%% ch02_background.tex — Chapter 2: Background and Related Work
%% HSP-Agile CCF-B Report, Ye Xujun, ECNU, July 2026
%% Staging rewrite (_rewrite/); do not edit live chapters/ from this file.

\chapter{Background and Related Work}
\label{ch:background}

Specification-conformance prevention draws on several mature ingredients:
ordered FSF contracts, LLM synthesis, CEGIS-style repair, SMT witnesses,
requirement-level pattern screens, and mutation-based scoring.
Prior systems typically assemble only a subset of these for
specification-conformance defects.
This chapter introduces the ingredients in concept order---SOFL and FSF
first, then the synthesis and verification mechanisms that operate on
them---and closes with a single positioning of \textsc{SgDP} relative to
adjacent lines of work.
Formal definitions are deferred to Chapter~\ref{ch:formalization}.

% ─────────────────────────────────────────────────────────────────────────────
\section{SOFL and Agile-SOFL}
\label{sec:bg:sofl}
% ─────────────────────────────────────────────────────────────────────────────

SOFL combines Z-style predicate logic, data-flow diagrams, and
object-oriented structuring under stepwise
refinement~\citep{liu1998sofl,back1998refinement}.
Its guarded process scenarios are the historical ancestor of the Functional
Specification Format (FSF) used throughout this report
(Section~\ref{sec:bg:fsf}).

\emph{Agile-SOFL}~\citep{li2022qrs_companion} places that rigour inside
sprint workflows: FSF contracts serve as design artefacts, oracles, and
acceptance criteria, with automatic conformance checking and SMT-directed
test generation before review~\citep{li2023iet_faultprevention,beck2004xp}.
HSP-Agile builds an LLM-driven synthesis--repair loop on that toolchain,
addressing the workflow-integration barrier identified by
\citet{woodcock2009formalsurvey}: formal artefacts help only when they
participate in everyday delivery, not when they remain specialist sidelines.

% ─────────────────────────────────────────────────────────────────────────────
\section{The Functional Specification Format (FSF)}
\label{sec:bg:fsf}
% ─────────────────────────────────────────────────────────────────────────────

FSF is the primary sprint artefact of
Agile-SOFL~\citep{li2022qrs_companion,liu2004soflbook}.
A task declares parameters, readable and writable externals, and an
\emph{ordered} sequence of scenario clauses.
Each clause pairs a Boolean guard with an assignment block; the final clause
may be an \texttt{others} catch-all and must appear last.
Guards are not a flat case statement: \emph{order is part of the contract}.

\subsection{First-match semantics and residual regions}
\label{sec:bg:fsf:semantics}

The defining rule is \emph{first-match}: the first satisfied guard determines
the output.
Two concrete failure modes follow from that rule; both recur under LLM
synthesis.

\paragraph{Primary vignette: catch-all residual.}
On the tier classifier of Chapter~\ref{ch:intro}
(Listing~\ref{lst:others-fsf}), named guards handle Reject and Pass bands;
everything else falls into \texttt{OTHERWISE} and must return
\texttt{Review}.
An LLM candidate that collapses the residual into \texttt{Pass} can satisfy
every busy-band unit test while shipping a wrong mid-band behaviour---false
acceptance in miniature.
The catch-all is thin: random suites drawn from Reject/Pass bands often miss
it entirely.

\paragraph{Secondary vignette: overlapping guards.}
On \code{classify\_signal}, input
$(\mathit{level},\mathit{threshold})=(-3,-10)$ satisfies both $S_1$
($\mathit{level}<0$) and $S_2$ ($\mathit{level}\ge\mathit{threshold}$), yet
FSF requires \texttt{"Error"} because $S_1$ has priority.
An implementation that tests the threshold first returns
\texttt{"Critical"}---a \emph{scenario-order violation} invisible to checkers
that treat guards as an unordered Boolean formula.

Chapter~\ref{ch:formalization} defines the oracle $g_t$ and the precedence
partition formally (Figure~\ref{fig:fsf-semantics}).
Unlike Dijkstra's guarded commands~\citep{dijkstra1975guarded}, FSF is
deterministic priority, not nondeterministic choice.

% ─────────────────────────────────────────────────────────────────────────────
\section{LLM-Assisted Code Generation}
\label{sec:bg:llm}
% ─────────────────────────────────────────────────────────────────────────────

Contest-style generation optimises for pass@k on hidden
tests~\citep{chen2021codex,li2022alphacode,openai2023gpt4}.
FSF synthesis needs something stricter: conformance to an ordered scenario
contract, including residual regions that coarse suites undersample.
Deployed assistants often produce \emph{semantically plausible} code that
fails at boundaries and multi-condition
logic~\citep{dakhel2023copilot}---precisely what FSF guards encode.

Without additional guidance, three failure modes recur
\citep{dakhel2023copilot,sun2024clover}:
(i)~\emph{catch-all / residual blindness}---hard-coded defaults that pass
busy-band tests (pattern RF07; Section~\ref{sec:bg:patterns});
(ii)~\emph{scenario-order blindness}---overlapping guards bias the model
toward a ``natural'' \texttt{if} order rather than mandated precedence;
(iii)~\emph{boundary hallucination}---confused relational operators at
scenario edges (e.g.\ $<$ vs.\ $\le$).
Test-feedback repair~\citep{chen2021codex} helps only when tests hit the
violated region; fixed suites often miss thin catch-alls and overlap slices.
HSP-Agile therefore synthesises region witnesses with Z3 rather than relying
on sampled unit tests alone.

% ─────────────────────────────────────────────────────────────────────────────
\section{Counterexample-Guided Synthesis (CEGIS)}
\label{sec:bg:cegis}
% ─────────────────────────────────────────────────────────────────────────────

Repair needs counterexamples that name the violated region, not only a failing
assertion.
CEGIS splits synthesis into a synthesiser and a verifier that returns
counterexamples until success or budget exhaustion
\citep{solarlezama2008sketching}.
LLM-augmented variants replace templates with free-form generation while
retaining formal checks~\citep{loughridge2024dafnybench,sun2024clover};
those pipelines target Dafny-style proof obligations rather than FSF
first-match regions or requirements-level pattern screens.

HSP-Agile instantiates the \textsc{SgDP} loop with an LLM synthesiser, a
Z3-based verifier over suite $\mathcal{D}(t)$, and a feedback aggregator that
packages counterexamples $(x^*, g_t(x^*), f_{c_k}(x^*))$ together with
high-severity pattern hits (Figure~\ref{fig:cegis-loop}).
Budget $K$ caps iterations; on exhaustion the highest-$\mathrm{Conf}$
candidate is returned as best-effort
($\arg\max_k \mathrm{Conf}(c_k,t)$)~\citep{barrett2018smt}.
Experimental budgets and mode definitions are given in
Chapter~\ref{ch:experiments}.

\begin{figure}[t]
  \centering
  \tikzinput{diagrams/tikz/cegis_loop}
  \caption{The HSP-Agile CEGIS loop.
           The LLM synthesiser produces candidate $c_k$ from FSF specification
           $\mathcal{S}_t$ and repair context.
           The conformance checker verifies $c_k$ against the Z3-generated test
           suite; failures become counterexamples.
           The feedback aggregator packages counterexamples and anti-patterns
           for the next iteration.
           If all tests pass and pattern guards are clear, $c_k$ is accepted.}
  \label{fig:cegis-loop}
\end{figure}

% ─────────────────────────────────────────────────────────────────────────────
\section{SMT Solving and Formal Verification}
\label{sec:bg:smt}
% ─────────────────────────────────────────────────────────────────────────────

SMT extends SAT with background theories; a model is a concrete test
input~\citep{barrett2018smt,demoura2008z3}.
For each scenario $s_i$, HSP-Agile solves the \emph{effective} guard
$\varphi_i = g_i(x) \wedge \bigwedge_{j<i}\neg g_j(x)$, guaranteeing at least
one witness per first-match region---including catch-alls encoded as the
negation of all prior guards.
That coverage is targeted rather than random: property-based testing
(e.g.\ QuickCheck) samples broad input families from executable
properties~\citep{claessen2000quickcheck,zhu1997specification}, but does not
by itself emit region-typed failure records for ordered FSF partitions.

The resulting oracle is weaker than full deductive verification
(e.g.\ Dafny~\citep{leino2010dafny}) yet requires no developer annotations and
fits Agile-SOFL sprint velocity.
Prevention guarantees remain bounded by suite coverage
(Chapter~\ref{ch:formalization}).

% ─────────────────────────────────────────────────────────────────────────────
\section{Requirement-Level Error Pattern Detection}
\label{sec:bg:patterns}
% ─────────────────────────────────────────────────────────────────────────────

Static analysers encode recurring structural bug
patterns~\citep{hovemeyer2004findbugs,cousot1977abstract}.
HSP-Agile's pattern guard lifts the idea to \emph{requirements-level}
violations: signatures that cannot conform to FSF regardless of tested inputs.
Catalogue RF01--RF14~\citep{li2023iet_faultprevention}
(Appendix~\ref{app:patterns}) includes missing preconditions, wrong defaults,
swapped comparisons, and unconditional success (RF07); high-severity hits
($\mathcal{P}^H$) block acceptance and enter CEGIS repair context even when
all formal tests pass.
Conjunctive release---witness agreement \emph{and} a clear pattern
screen---is the Accept predicate developed in later chapters.

% ─────────────────────────────────────────────────────────────────────────────
\section{Mutation Testing for Specification Assessment}
\label{sec:bg:mutation}
% ─────────────────────────────────────────────────────────────────────────────

Mutation testing injects syntactic faults; survivors indicate inadequate
coverage~\citep{demillo1978mutation,offutt2001mutation2000,ammann2016testing}.
HSP-Agile uses two operator families (details in Chapter~\ref{ch:experiments}):
specification-level SNO, ORO, SPO; implementation-level ICO, WRO, BCO.
Prevention quality is scored as \emph{Prevention Detection Rate} (PDR: fraction
of injected faults correctly rejected) and \emph{False Accept Rate} (FAR:
escape rate among faulty mutants), both over the faulty-mutant set
(Definitions~\ref{def:pdr}--\ref{def:far}).
These metrics complement Conf.\ / pass-rate: a candidate may clear a fixed suite
yet still accept faulty mutants that encode residual or ordering defects.

% ─────────────────────────────────────────────────────────────────────────────
\section{Related Work}
\label{sec:bg:related}
% ─────────────────────────────────────────────────────────────────────────────

What is missing is not another synthesiser, but a prevention loop that treats
ordered guards as first-class evidence and scores release by prevention, not
only pass@k.
Table~\ref{tab:related-work} summarises the comparison on the axes that matter
for this report.

Successful formal-method deployments automate narrow, high-value layers inside
existing workflows~\citep{woodcock2009formalsurvey,pohl2010requirements};
Agile-SOFL follows that prescription, and HSP-Agile automates synthesis and
repair without manual annotation.
Empirical SE guidance on construct and external validity
\cite{wohlin2012experimentation,ammann2016testing} shapes how we separate
mechanism claims (typed feedback under E6) from deployment heuristics (when to
escalate from a lighter loop to the full Accept gate).

\subsection{LLM synthesis, agents, and automated program repair}
Contest-style generation optimises for pass@k on hidden
tests~\citep{chen2021codex,li2022alphacode,openai2023gpt4,austin2021mbpp,lu2021codexglue}.
Repository agents (OpenHands, SWE-agent, Agentless, AutoCodeRover) and
multi-agent coders target GitHub-scale
tasks~\citep{wang2024opendevin,yang2024sweagent,xia2025agentless,zhang2024autocoderover,jimenez2024swebench,huang2023agentcoder}
without a first-match specification oracle $g_t$ or PDR/FAR protocols.
Classical and LLM-based APR (GenProg, TBar, ChatRepair, RepairAgent, and
related studies) patch failing tests or buggy
snippets~\citep{legoues2019apr,liu2019tbar,xia2024chatrepair,bouzenia2024repairagent,jiang2023impact,fan2023llmrepair,sobania2023chatgptapr,wei2023copiloteval,olausson2024selfrepair,yasunaga2021breakitfix};
our loop instead repairs against \emph{ordered-guard} witnesses derived from
SpecIR and scores prevention, not only pass@k.
Verbal self-repair (Self-Refine, Self-Debug, Reflexion) raises generation
success via critique or memory~\citep{madaan2023selfrefine,chen2024selfdebug,shinn2023reflexion}
but does not emit region-typed failure records keyed to scenario identity.

\subsection{LLM-as-judge and verbal structured feedback}
A natural alternative to typed IR is to ask another LLM to critique the
candidate in natural language---LLM-as-judge / self-refine style
feedback~\citep{madaan2023selfrefine,chen2024selfdebug,shinn2023reflexion}.
Those loops can raise generation success, but the critique is not bound to an
executable first-match oracle $g_t$: the judge may invent a plausible story,
omit the catch-all scenario index, or disagree with $\Phi_i$ witnesses.
Semantic Feedback IR is the opposite design choice: structure is
\emph{derived from} SMT witnesses and scenario identity, then rendered to
text---not elicited as free-form opinion.
Baselines B3--B5 instantiate verbal/self-critique styles under equal~$K$ on
BENCH-120 (Chapter~\ref{ch:experiments}); they are related feedback channels,
not substitutes for oracle-grounded region records.

\subsection{Oracles, property-based testing, and symbolic suites}
The oracle problem remains central whenever tests are the only
gate~\citep{barr2015oracle,hierons2012oracle,zhu1997specification}.
Property-based testing (QuickCheck, Hypothesis) samples broad input families
from executable properties~\citep{claessen2000quickcheck,maciver2019hypothesis};
search-based and symbolic generators (EvoSuite, KLEE, Pex) raise structural
coverage~\citep{fraser2011evosuite,cadar2008klee,tillmann2008pex,anand2013symbolic}.
HSP-Agile complements those lines with solver-directed witnesses for
\emph{first-match} regions $\Phi_i$, then packages failures as typed repair
records rather than coverage maximisation alone.
Mutation testing surveys and operators inform our PDR/FAR
protocol~\citep{jia2011mutation,demillo1978mutation,ammann2016testing}.

\subsection{Contracts, proof-oriented LLM loops, and adjacent notations}
Design-by-Contract and embedded contract languages provide runtime
oracles~\citep{meyer1997oosc,fahndrich2010codecontracts} but do not encode FSF
precedence partitions.
DafnyBench and Clover-style verifier-in-the-loop systems target proof
obligations~\citep{loughridge2024dafnybench,sun2024clover,first2023baldur,kamath2023invariants,moura2021lean4};
they are fair references for closed-loop repair, not drop-in FSF comparators.
Alloy and TLA+ encode structured
constraints~\citep{jackson2002alloy,lamport2002tla}; SpecIR adapters are the
intended bridge when guards are first-match and SMT-encodable.
Prior SOFL fault-prevention work assumed human-written
implementations~\citep{li2022qrs_companion,li2023iet_faultprevention}; LLM
synthesis introduces residual and ordering failure modes that neither
human-process prevention nor test-only post-processing alone addresses.

Closest LLM+verifier systems differ on three axes: whether witnesses respect
first-match order, whether repair carries scenario structure, and whether
prevention (not only pass@k) is measured.
Our controlled baselines isolate those ingredients:
B1 (one-shot LLM), B2 (test-feedback without SMT witnesses), B6
(VerifierLoop-FSF: SMT counterexamples without typed Semantic Feedback IR or
a hard pattern screen), and M (full \textsc{SgDP} with IR and Accept gate).
Primary ranking and mechanism results appear in
Chapter~\ref{ch:results}; the comparison axes below are design distinctions,
not a substitute for those tables.

\Needspace{14\baselineskip}
\begin{inlinetable}{Structured comparison of related approaches to LLM-assisted
           specification-conformant code generation.
           Checkmarks (\checkmark) indicate full support; dashes (--) absence;
           ``partial'' indicates limited support.
           Columns emphasise the axes that matter for this paper:
           first-match witnesses, typed repair IR, prevention metrics, and
           release gate.}
  \label{tab:related-work}
  \label{tab:related-positioning}% retained alias: merged into this single table
  \footnotesize
  \renewcommand{\arraystretch}{1.12}
  \begin{tabular*}{\linewidth}{@{\extracolsep{\fill}}l c c c c c@{}}
    \toprule
    \textbf{Approach}
      & \textbf{Ordered $\Phi_i$}
      & \textbf{Typed IR}
      & \textbf{Prev.}
      & \textbf{Gate}
      & \textbf{Primary metric} \\
    \midrule
    DafnyBench \citep{loughridge2024dafnybench}
      & --
      & --
      & --
      & proof
      & pass@k \\
    Sun et al.\ (Clover) \citep{sun2024clover}
      & --
      & partial
      & --
      & proof
      & correctness \\
    LLM-APR / agents \citep{xia2024chatrepair,xia2025agentless}
      & --
      & --
      & --
      & tests
      & pass@k \\
    Self-Refine / Self-Debug / Reflexion
      \citep{madaan2023selfrefine,chen2024selfdebug,shinn2023reflexion}
      & --
      & --
      & --
      & verbal
      & pass@k \\
    Standard LLM (B1)
      & --
      & --
      & --
      & --
      & pass rate \\
    Test-feedback (B2)
      & --
      & --
      & --
      & tests
      & Conf.\ / pass \\
    B6 VerifierLoop-FSF
      & \checkmark
      & --
      & --
      & cex
      & Conf. \\
    \midrule
    \textbf{HSP-Agile (ours)}
      & \checkmark
      & \checkmark
      & \checkmark
      & Accept
      & \textbf{Conf, PDR, FAR} \\
    \bottomrule
  \end{tabular*}%
\end{inlinetable}

\noindent
Relative to these lines of work, \textsc{SgDP}'s increment is the
\emph{combination} of FSF first-match witnesses, specification-indexed
Semantic Feedback IR, and conjunctive $\mathrm{Accept}$ under PDR/FAR
evaluation~\citep{chen2021codex,loughridge2024dafnybench,sun2024clover,legoues2019apr}.
Versus B6, the design delta is typed IR plus the Accept/pattern gate---SMT
counterexamples alone do not constitute the full loop.
Versus APR and verbal self-debug, the delta is SpecIR-region witnesses and
prevention metrics rather than pass@k alone.
Versus Clover-style verifier loops, success here is witness agreement under
first-match partitions plus a release predicate, not theorem discharge.
Property-based testing and symbolic suites raise coverage without emitting
region-typed repair records~\citep{claessen2000quickcheck,cadar2008klee};
metamorphic and oracle-problem surveys frame why executable SpecIR oracles
matter when tests alone undersample thin guards~\citep{barr2015oracle,hierons2012oracle}.

% ─────────────────────────────────────────────────────────────────────────────
\section{Gap Summary}
\label{sec:bg:gap}
% ─────────────────────────────────────────────────────────────────────────────

Table~\ref{tab:related-work} leaves open three design commitments that
Chapter~\ref{ch:formalization} makes precise---with SpecIR as the shared
engineering hinge:

\begin{enumerate}
  \item[\textbf{Premise.}] Prevention must depend on \emph{semantic regions}
        (first-match partitions), not surface syntax---realised by SpecIR.
  \item[\textbf{C2.}] Repair needs \emph{specification-indexed structured
        feedback}, not raw pass/fail bits---realised by the Semantic Feedback
        IR (controlled contrast: scenario-indexed IR vs.\ test-only dumps).
  \item[\textbf{C3.}] Release must be \emph{conjunctive}: every witness matches
        \emph{and} no high-severity structural pattern remains
        ($\mathrm{Accept} \Leftrightarrow \mathrm{Conf}{=}1 \wedge
        \mathrm{Screen}{=}\mathrm{pass}$).
\end{enumerate}

Deployment guidance (C4) then chooses when to pay for the full loop versus a
lighter test-feedback baseline (B2): escalate when Accept/FAR headroom or
residual-region risk warrants it.
The catch-all vignette of Chapter~\ref{ch:intro} remains the conceptual
anchor for those claims; overlapping-guard tasks are a secondary stress on
ordering.
Empirical ranking of baselines, equal-$K$ Conf.\ comparisons, and prevention
(PDR/FAR) appear in Chapters~\ref{ch:experiments}--\ref{ch:results}.

